\documentclass[11pt,a4paper]{article}
\usepackage[margin=25mm]{geometry}
\usepackage{amsmath,amssymb,amsthm}
\newtheorem{proposition}{Proposition}
\title{A boundary case for the rank-four map construction}
\author{Human Compression project: L023 working audit}\date{}
\begin{document}\maketitle
This short, unreviewed calculation concerns Construction~6, Remark~9 and
the general class correspondence on page~11 of Leemans--Stokes,
arXiv:2208.08215v2. It does not disprove their Class~III Theorem~10.
The counterexample is a reflexible map with two vertices, two parallel
edges and two digonal faces on the sphere; it has eight flags.

Let $G=\mathbb F_2^3$ additively, with basis $a,b,c$, and take the
three marked involutions to be $a,b,c$ in that order. A flag of this map
can be indexed by the choice of its vertex, edge and face. Changing any
one choice gives the corresponding flag adjacency. These three commuting
involutions act regularly on the eight flags. In particular this is an
actual map with faithful incidence action, as well as a quotient of the
source's universal map group. All three periods are two.
Every linear automorphism fixing $b$ and permuting the three nonzero
vectors in $\langle a,c\rangle$ is a marked Wilson symmetry. There
are all six permutations, so the map belongs to Class~IV.

Construction~6 gives the four subgroups
\[
 H_0=\langle a,b\rangle,\quad H_1=\langle a,c\rangle,
 \quad H_2=\langle b,c\rangle,\quad H_3=\langle b,a+c\rangle.
\]
They are distinct planes in $G$.

\begin{proposition}
The resulting rank-four coset geometry is the complete four-partite
incidence structure with two objects of each type. It is firm, and
\[
 \operatorname{Aut}(\Gamma)\cong C_2^4,
 \qquad \operatorname{Cor}(\Gamma)\cong C_2^4\rtimes S_4.
\]
Thus $\operatorname{Cor}(\Gamma)/\operatorname{Aut}(\Gamma)\cong S_4$,
rather than the source's asserted $S_3$ for Class~IV.
\end{proposition}
\begin{proof}
Each subgroup has index two. Distinct planes have sum $G$, so every coset
of one meets every coset of another: for cosets $x+H_i,y+H_j$, express
$y-x=u-v$ with $u\in H_i,v\in H_j$ to obtain the common point
$x+u=y+v$. Thus all cross-type pairs are incident. Every maximal flag
therefore chooses one object from each of the four pairs; there are
sixteen chambers. A flag missing one type has exactly two completions,
so this geometry is firm.

Independent swaps in the four pairs give every type-preserving
permutation, and all preserve incidence. This is $C_2^4$.
Any permutation of the four types, with independent bijections on their
pairs, is a correlation, and these exhaust correlations by definition.
Their group is $C_2^4\rtimes S_4$, proving the assertions.
\end{proof}
The face--Petrie residue consists here of two vertices and two edges,
with every cross pair incident. It is a four-cycle, not a path on three
vertices and two edges. Thus the unrestricted residue description used
in source Lemma~8 and Remark~9 fails. The numerical conclusion of
Lemma~8 happens to hold for this example: the action of $G$ on chamber
labels $(x_0,x_1,x_2,x_3)\in\mathbb F_2^4$ has image
\[
 \{(C,B,A,A+C):A,B,C\in\mathbb F_2\},
\]
so has exactly two chamber orbits. No counterexample to that numerical
claim is inferred from this calculation.

This is an exact written counterexample to the unrestricted nonfirmness
and Class~IV correlation-quotient assertions, with no finite computation
as premise. Restricting a proposed replacement theorem to another class
or to stronger incidence hypotheses would require a separate proof.
The general flag-extension claim and the Class~III no-duality conclusion
remain separate recovery obligations.
\end{document}
